\documentclass[11pt]{article}

% -------------------- Packages --------------------
\usepackage[a4paper,margin=1in]{geometry}
\usepackage{hyperref}
\usepackage{enumitem}
\usepackage{titlesec}
\usepackage{fancyhdr}
\usepackage{xcolor}
\usepackage{listings}
\usepackage{courier}

% -------------------- Hyperref Setup --------------------
\hypersetup{
  colorlinks=true,
  linkcolor=blue,
  urlcolor=blue,
  citecolor=blue
}

% -------------------- Header/Footer --------------------
\pagestyle{fancy}
\fancyhf{}
\lhead{MCV Lab}
\rhead{API Design \& Scaling}
\cfoot{\thepage}

% -------------------- Section Formatting --------------------
\titleformat{\section}{\Large\bfseries}{\thesection}{0.8em}{}
\titleformat{\subsection}{\normalsize\bfseries}{\thesubsection}{0.8em}{}
\titleformat{\subsubsection}{\normalsize\bfseries}{\thesubsubsection}{0.8em}{}

\titlespacing*{\section}{0pt}{1.5em}{0.8em}
\titlespacing*{\subsection}{0pt}{1.0em}{0.5em}
\titlespacing*{\subsubsection}{0pt}{0.8em}{0.4em}

% -------------------- Global List Formatting --------------------
\setlist[itemize]{
  label={-},
  leftmargin=2em,
  labelsep=0.6em,
  itemsep=0.2em
}
\setlist[itemize,2]{itemsep=0.15em}
\setlist[enumerate]{leftmargin=2em,itemsep=0.15em}

% -------------------- Listings / Code Style --------------------
\definecolor{codegray}{rgb}{0.95,0.95,0.95}
\definecolor{commentgreen}{rgb}{0.0,0.5,0.0}

\lstset{
  basicstyle=\footnotesize\ttfamily,
  backgroundcolor=\color{codegray},
  frame=single,
  breaklines=true,
  showstringspaces=false,
  keywordstyle=\color{blue},
  commentstyle=\color{commentgreen},
  stringstyle=\color{red},
  columns=fullflexible,
  keepspaces=true,
  tabsize=2,
  xleftmargin=0.02\textwidth,
  xrightmargin=0.02\textwidth
}

% -------------------- Document Metadata --------------------
\title{\textbf{Medical Computer Vision Lab \\ \small{
Designing and Scaling APIs/Endpoints (Flask, Django, FastAPI): Lab Standard}}}
\author{Maintainer: Ashery Mbilinyi}
\date{Version 1.0 : \today}

\begin{document}
\maketitle

\tableofcontents
\newpage

\section{Purpose}
This document defines MCV Lab standards for \textbf{designing} and \textbf{scaling} APIs/endpoints for ML systems.
The goal is to ensure our services are:
\begin{itemize}
  \item \textbf{Correct}: consistent request/response contracts and validation
  \item \textbf{Reliable}: robust error handling, health checks, and observability
  \item \textbf{Maintainable}: clear project structure and versioning
  \item \textbf{Scalable}: supports increased traffic and safe deployment patterns
\end{itemize}

\section{Core Concepts}

\subsection{API and Endpoint}
\begin{itemize}
  \item \textbf{API}: the set of routes exposed by a service.
  \item \textbf{Endpoint}: a specific route (e.g., \texttt{POST /predict}, \texttt{GET /health}).
\end{itemize}

\subsection{Common ML Service Endpoints (Recommended)}
All ML inference services should implement:
\begin{itemize}
  \item \texttt{GET /health}: process is alive (liveness)
  \item \texttt{GET /ready}: model is loaded and ready (readiness)
  \item \texttt{GET /version}: service version and model version
  \item \texttt{POST /predict}: inference on a single request
  \item \texttt{POST /batch\_predict} (optional): batch inference for throughput
\end{itemize}

\subsection{Contract First: Request/Response Schema}
Each endpoint must have:
\begin{itemize}
  \item a well-defined request schema (types, required fields, ranges)
  \item a well-defined response schema
  \item clear error responses with HTTP status codes
\end{itemize}

\section{Framework Selection Guide}

\subsection{Flask}
\textbf{Best for:} simple prototypes and lightweight services.\\
\textbf{Tradeoff:} minimal structure; you must add validation and documentation manually.

\subsection{Django (with Django REST Framework)}
\textbf{Best for:} full platforms with authentication, database models, admin UI, and complex workflows.\\
\textbf{Tradeoff:} heavier than needed for pure inference microservices.

\subsection{FastAPI (Recommended Default)}
\textbf{Best for:} modern inference microservices with typed schemas and automatic docs.\\
\textbf{Why we like it:} Pydantic validation + OpenAPI docs + clean contracts.

\section{Scaling Concepts (What ``Scaling an API'' Means)}

\subsection{Concurrency}
Ability to serve multiple requests at once.
Common approaches:
\begin{itemize}
  \item multiple worker processes (e.g., Gunicorn)
  \item async I/O (FastAPI) for network/disk-bound operations
\end{itemize}
\textbf{Important:} ML inference is typically CPU/GPU-bound; async does not magically speed up inference.

\subsection{Horizontal Scaling}
Run multiple replicas of the API service and load-balance across them.

\subsection{Batching and Throughput (Critical for GPUs)}
Batch inference improves throughput by using GPU efficiently.
If latency constraints allow, prefer batching.

\subsection{Load Balancing}
Use a reverse proxy (Nginx) or platform load balancer (Kubernetes Service/Ingress).

\section{Lab Policy}

\subsection{Minimum Requirements for Any API Service}
\begin{itemize}
  \item \texttt{/health}, \texttt{/ready}, \texttt{/version}, \texttt{/predict} implemented
  \item request validation present (FastAPI/Pydantic strongly preferred)
  \item structured logging (JSON or consistent text format)
  \item deterministic model versioning (model id/hash + git commit in \texttt{/version})
  \item containerized deployment recipe (Dockerfile)
\end{itemize}

\subsection{Recommended Requirements for Production-Style Services}
\begin{itemize}
  \item rate limiting (basic protection)
  \item authentication/authorization if exposing beyond local network
  \item metrics endpoint (Prometheus) or equivalent
  \item safe rollout and rollback strategy
\end{itemize}

\section{Coding Examples}

\subsection{FastAPI Example (Recommended Inference Service)}
\subsubsection{app.py}
\begin{lstlisting}[language=Python]
from __future__ import annotations

import os
from typing import List, Optional

from fastapi import FastAPI, HTTPException
from pydantic import BaseModel, Field

APP_VERSION = os.getenv("APP_VERSION", "dev")
MODEL_VERSION = os.getenv("MODEL_VERSION", "unknown")

app = FastAPI(title="MCV Inference API", version=APP_VERSION)

# ---- Schema (Contract) ----
class PredictRequest(BaseModel):
    # Example: a vector input. Replace with your real schema.
    x: List[float] = Field(..., min_length=1, description="Input feature vector.")
    return_probs: bool = False

class PredictResponse(BaseModel):
    y: int
    probs: Optional[List[float]] = None
    model_version: str

# ---- Model loading (placeholder) ----
MODEL_LOADED = True

def dummy_model_predict(x: List[float]) -> tuple[int, List[float]]:
    # Replace with real model inference
    # Example rule:
    score = sum(x)
    y = 1 if score > 0 else 0
    probs = [1.0 - 0.2, 0.2] if y == 0 else [0.2, 0.8]
    return y, probs

# ---- Health endpoints ----
@app.get("/health")
def health():
    return {"status": "ok"}

@app.get("/ready")
def ready():
    if not MODEL_LOADED:
        raise HTTPException(status_code=503, detail="Model not loaded")
    return {"status": "ready"}

@app.get("/version")
def version():
    return {"app_version": APP_VERSION, "model_version": MODEL_VERSION}

# ---- Inference ----
@app.post("/predict", response_model=PredictResponse)
def predict(req: PredictRequest):
    if not MODEL_LOADED:
        raise HTTPException(status_code=503, detail="Model not loaded")

    y, probs = dummy_model_predict(req.x)

    return PredictResponse(
        y=y,
        probs=probs if req.return_probs else None,
        model_version=MODEL_VERSION,
    )
\end{lstlisting}

\subsubsection{Run locally}
\begin{lstlisting}[language=bash]
pip install fastapi uvicorn pydantic
uvicorn app:app --host 0.0.0.0 --port 8000
\end{lstlisting}

\subsubsection{Test}
\begin{lstlisting}[language=bash]
curl http://localhost:8000/health
curl http://localhost:8000/version

curl -X POST http://localhost:8000/predict \
  -H "Content-Type: application/json" \
  -d '{"x":[1.0,-0.5,0.2], "return_probs": true}'
\end{lstlisting}

\subsection{Flask Example (Minimal Prototype)}
\begin{lstlisting}[language=Python]
from flask import Flask, request, jsonify

app = Flask(__name__)
MODEL_LOADED = True

@app.get("/health")
def health():
    return jsonify({"status": "ok"})

@app.get("/ready")
def ready():
    if not MODEL_LOADED:
        return jsonify({"status": "not_ready"}), 503
    return jsonify({"status": "ready"})

@app.post("/predict")
def predict():
    payload = request.get_json(force=True)
    x = payload.get("x")
    if x is None:
        return jsonify({"error": "Missing field: x"}), 400
    y = 1 if sum(x) > 0 else 0
    return jsonify({"y": y})
\end{lstlisting}

\subsection{Django REST Framework (DRF) Sketch}
Use DRF if you need user accounts, permissions, and database-backed workflows.
\begin{lstlisting}[language=Python]
# serializers.py
from rest_framework import serializers

class PredictSerializer(serializers.Serializer):
    x = serializers.ListField(
        child=serializers.FloatField(), allow_empty=False
    )

# views.py
from rest_framework.views import APIView
from rest_framework.response import Response
from rest_framework import status
from .serializers import PredictSerializer

class PredictView(APIView):
    def post(self, request):
        s = PredictSerializer(data=request.data)
        s.is_valid(raise_exception=True)
        x = s.validated_data["x"]
        y = 1 if sum(x) > 0 else 0
        return Response({"y": y}, status=status.HTTP_200_OK)
\end{lstlisting}

\section{Scaling Recipes}

\subsection{Production Server (FastAPI)}
For production, do not rely on a single dev server. Use multiple workers.
\begin{lstlisting}[language=bash]
pip install "uvicorn[standard]" gunicorn

gunicorn -k uvicorn.workers.UvicornWorker \
  -w 2 \
  -b 0.0.0.0:8000 \
  app:app
\end{lstlisting}

\subsection{Reverse Proxy (Nginx) --- Minimal Concept}
Nginx can terminate TLS and load balance across multiple app workers/replicas.

\subsection{Kubernetes (High Level)}
Kubernetes scales by running multiple replicas:
\begin{itemize}
  \item \textbf{Deployment}: defines desired replicas
  \item \textbf{Service}: load-balances traffic across replicas
  \item \textbf{HPA}: autoscaling based on CPU/custom metrics
\end{itemize}

\section{Common Pitfalls and Guidance}

\subsection{Async Misconceptions}
Async helps with I/O, not with GPU-bound inference. For GPU inference:
\begin{itemize}
  \item use batching where possible
  \item consider a queue-based design for heavy requests
  \item avoid too many workers that compete for the same GPU memory
\end{itemize}

\subsection{Model Loading}
\begin{itemize}
  \item load model once at startup (not per request)
  \item warm up the model (optional) to stabilize latency
  \item expose readiness only when model is loaded
\end{itemize}

\subsection{Versioning}
Always expose:
\begin{itemize}
  \item git commit hash (app version)
  \item model artifact id/hash (model version)
\end{itemize}

\section{Checklist (Code Review)}
Before merging an API service:
\begin{itemize}
  \item endpoints: \texttt{/health}, \texttt{/ready}, \texttt{/version}, \texttt{/predict}
  \item request/response schema is documented and validated
  \item errors return correct HTTP codes (400 vs 503 vs 500)
  \item model loads once and readiness reflects true state
  \item production run command documented (Gunicorn/Uvicorn)
  \item Dockerfile present
\end{itemize}

\end{document}
